\begin{figure}[htbp]
	\centering        
	\begin{tikzpicture}[>=Stealth, line width=0.8pt]
		
		% ═══════════════════════════════════════════════════════════
		% PANNELLO (a) — Reazione del carrello
		% ═══════════════════════════════════════════════════════════
		\begin{scope}[xshift=0cm]
			
			% ── Parametri ──
			\def\xA{0.8}
			\def\yA{0.0}
			\def\tw{0.30}
			\def\thR{0.40}
			\def\rn{0.10}
			\def\bw{0.18}
			\def\gw{0.45}
			
			% ── Asse del carrello: retta verticale grigia continua ──
			\draw[gray!60, line width=0.6pt] 
			(\xA, \yA-\thR-0.5) -- (\xA, 3.85);
			
			% ── Corpo di fantasia ──
			\draw[line width=1pt] 
			plot[smooth cycle, tension=0.7] coordinates 
			{(0.8, 0.0) (1.6, 0.15) (2.2, 0.7) (2.0, 1.7) 
				(1.0, 2.0) (-0.1, 1.7) (-0.6, 1.0) (-0.3, 0.3)};
			\node at (1.5, 1.5) {$\mathcal{C}$};
			
			% ── Simbolo del carrello ──
			\draw (\xA, \yA) -- ++(-\tw, -\thR) -- ++(2*\tw, 0) -- cycle;
			\draw (\xA-\gw, \yA-\thR) -- (\xA+\gw, \yA-\thR);
			\draw (\xA-\bw, \yA-\thR-0.13) circle (3pt);
			\draw (\xA+\bw, \yA-\thR-0.13) circle (3pt);
			\draw (\xA-\gw, \yA-\thR-0.26) -- (\xA+\gw, \yA-\thR-0.26);
			\foreach \x in {-0.36,-0.18,0,0.18,0.36}
			\draw[thin] (\xA+\x, \yA-\thR-0.26) -- ++(-0.10, -0.13);
			
			% ── Semicerchio in A: rotazione consentita ──
			\draw[fill=white] 
			(\xA-\rn, \yA) arc(180:360:\rn) -- cycle;
			
			% ── REAZIONE: forza R(A) a, in rosso, da A verso l'alto ──
			\draw[->, line width=1.2pt] 
			(\xA, \yA+0.05) -- (\xA, \yA+1.05);
			\node[anchor=west, font=\small] 
			at (\xA, \yA+0.65) {$R(A)\,\ab$};
			
			% ── Etichetta A (a destra, fuori dal triangolo) ──
			\node[anchor=west, font=\small] at (\xA+0.18, \yA-0.08) {$A$};
			
			% ── Asse a sopra il corpo ──
			\draw[->, thick] (\xA, 2.10) -- (\xA, 2.85);
			\node[right=2pt] at (\xA, 2.85) {$\ab$};
			
			% ── Etichetta verticale "asse del carrello" ──
			\node[gray, rotate=90, anchor=south, font=\footnotesize] 
			at (\xA-0.12, 2.45) {asse del carrello};
			
			% ── Etichetta pannello ──
			\node[anchor=north] at (\xA+2, \yA-\thR) {$(a)$};
			
		\end{scope}
		
		
		% ═══════════════════════════════════════════════════════════
		% PANNELLO (b) — Reazione della biella
		% ═══════════════════════════════════════════════════════════
		\begin{scope}[xshift=5.5cm]
			
			% ── Parametri ──
			\def\xA{0.8}
			\def\yA{1.5}
			\def\yG{0.0}
			\def\tw{0.30}
			\def\thA{0.40}
			\def\rn{0.10}
			\def\gw{0.50}
			
			% ── Asse della biella ──
			\draw[gray!60, line width=0.6pt] 
			(\xA, \yG-\thA-0.5) -- (\xA, 4.50);
			
			% ── Corpo di fantasia ──
			\draw[line width=1pt] 
			plot[smooth cycle, tension=0.7] coordinates 
			{(0.8, 1.5) (1.6, 1.65) (2.2, 2.2) (2.0, 3.2) 
				(1.0, 3.5) (-0.1, 3.2) (-0.6, 2.5) (-0.3, 1.8)};
			\node at (1.5, 3.0) {$\mathcal{C}$};
			
			% ── Asta della biella ──
			\draw[line width=0.8pt] (\xA, \yA) -- (\xA, \yG);
			
			% ── Cerniera al suolo ──
			\draw (\xA, \yG) -- ++(-\tw, -\thA) -- ++(2*\tw, 0) -- cycle;
			\draw (\xA-\gw, \yG-\thA) -- (\xA+\gw, \yG-\thA);
			\foreach \x in {-0.36,-0.18,0,0.18,0.36}
			\draw[thin] (\xA+\x, \yG-\thA) -- ++(-0.10, -0.13);
			
			% ── Anellini: semicerchio in A, perno pieno al suolo ──
			\draw[fill=white] 
			(\xA-\rn, \yA) arc(180:360:\rn) -- cycle;
			\draw[fill=white] (\xA, \yG) circle (3.5pt);
			
			% ── REAZIONE: forza R(A) a, in rosso, da A verso l'alto ──
			\draw[->,  line width=1.2pt] 
			(\xA, \yA+0.05) -- (\xA, \yA+1.05);
			\node[anchor=west, font=\small] 
			at (\xA, \yA+0.65) {$R(A)\,\ab$};
			
			% ── Etichetta A (a sinistra dell'asta) ──
			\node[anchor=east, font=\small] at (\xA, \yA-0.24) {$A$};
			
			% ── Asse a sopra il corpo ──
			\draw[->, thick] (\xA, 3.60) -- (\xA, 4.30);
			\node[right=2pt] at (\xA, 4.30) {$\ab$};
			
			% ── Etichetta verticale "asse della biella" ──
			\node[gray, rotate=90, anchor=south, font=\footnotesize] 
			at (\xA-0.12, 3.40) {asse della biella};
			
			% ── Etichetta pannello ──
			\node[anchor=north] at (\xA+2, \yG-\thA) {$(b)$};
			
		\end{scope}
		
	\end{tikzpicture}
	\caption{Prestazione statica del carrello (a) e della biella
		(b): la reazione vincolare si riduce alla sola forza
		$R(A)\,\ab$ diretta lungo l'asse, in dualità con la
		componente di spostamento bloccata; le componenti reattive
		nulle (coppia e forza in $\ab^\perp$) non sono
		rappresentate, in accordo con la molteplicità uno. Nella
		biella, la forza è trasmessa al corpo $\mathcal{C}$
		dall'asta lungo il proprio asse.}
	\label{fig:carrello_stat}
\end{figure}
